% !TEX TS-program = xelatex
% !TEX encoding = UTF-8 Unicode
% !Mode:: "TeX:UTF-8"

\documentclass{resume}
% \usepackage{zh_CN-Adobefonts_external} % Simplified Chinese Support using external fonts (./fonts/zh_CN-Adobe/)
%\usepackage{zh_CN-Adobefonts_internal} % Simplified Chinese Support using system fonts
\usepackage{linespacing_fix} % disable extra space before next section
\usepackage{cite}
\usepackage{ragged2e}
\usepackage{array}
\usepackage{setspace}
\usepackage{xcolor}
\usepackage{arydshln}
\usepackage{xeCJK}
% \usepackage{ulem}
% \usepackage{enumerate}
\usepackage{hyperref}
\usepackage{enumitem}
% \setlength{\parskip}{5em}
\hypersetup{
    colorlinks=true,            %链接颜色
    % linkcolor=blue,             %内部链接
    % filecolor=magenta,          %本地文档
    urlcolor=blue,              %网址链接
    % pdftitle={CV of Yang JS},
    % pdfpagemode=FullScreen,
    }

\begin{document}
% \pagenumbering{gobble} % suppress displaying page number
\pagenumbering{arabic}

% \section{\faAddressCard\,  基本信息}
% \begin{center}
  % \begin{minipage}[t]{0.75\textwidth}
    \vspace{0pt}
    % \hspace{-6pt}\name{Jia-Shu Yang \, \Large{Ph.D. Candidate}} 
    \hspace{-6pt}\name{Jia-Shu Yang \, \Large{Ph.D.}}
    \vspace{8pt}
    % \basicInfo{
        % \email{yang\_jiashu@outlook.com} \textperiodcentered\ 
        % \phone{+86 18717811089} \textperiodcentered\ }
        % \leftline{{\faEnvelope}\ jiashuyang@tongji.edu.cn \ \quad
        % {\faPhone}\ +86 18717811089 \ }
        
        % \leftline{{\faAddressBook}\ 1239 Siping Road, Shanghai 200092, China}
        % \leftline{{\faResearchgate}\ https://www.researchgate.net/profile/Jia-Shu-Yang}  
        % }      \vspace{6pt}
        
    \begin{spacing}{1.5}
    % \hspace{-7pt}
        {\faEnvelope}\; yang\_jiashu@outlook.com \qquad \qquad {\faPhone}\; +86 18717811089 \\
        {\faMapMarker*}\; Xi'an University of Architecture and Technology, 13 Yanta Road, Xi'an, Shaanxi 710055, China. \\
        {\faUserTie}\; https://www.researchgate.net/profile/Jia-Shu-Yang
        \vspace{8pt}
    \end{spacing}
  % \end{minipage}
  % \begin{minipage}[t]{0.24\textwidth}
    % \vspace{0pt}
    % \centering
    % \rightline{\includegraphics[width=0.6\textwidth]{resume/2021_pic11.jpg}}
  % \end{minipage}
% \end{center}

\vspace{-12pt}
\section{\faUniversity\ WORK EXPERIENCE}
\subsection{$\bullet$\;\;\textbf{Jul. 2023 -- present, Xi'an University of Architecture and Technology, Xi'an, China}}
\begin{spacing}{1.1}
 \quad\  Faculty-oriented postdoctoral fellow, School of Civil Engineering \par
\end{spacing}

% \section{\textcolor[rgb]{0,0.44,0.75}{\faGraduationCap\  教育背景}}
\section{\faGraduationCap\  EDUCATION}
%\subsection{$\bullet$\;\;\textbf{Sep. 2016 -- Present, Tongji University, Shanghai, China}}
\subsection{$\bullet$\;\;\textbf{Sep. 2016 -- Feb. 2023, Tongji University, Shanghai, China}}
\begin{spacing}{1.1}
\quad\ Ph.D. of Engineering, Department of Structural Engineering, College of Civil Engineering \par
\quad\ Supervisor: Prof. Jian-Bing Chen \par
%\quad\ Dissertation: Methods for seismic reliability-based structural optimization
\quad\ Dissertation: Seismic-Reliability-Based Structural Optimization by Probability Density Evolution Method
\end{spacing}
\subsection{$\bullet$\;\;\textbf{Jul. 2021 -- Jul. 2022, Leibniz Universität Hannover, Hannover, Germany}}
\begin{spacing}{1.1}
 \quad\ Visiting Ph.D. Student, Institute for Risk and Reliability \par
\quad\ Co-Supervisor: Prof. Michael Beer\
\end{spacing}
\subsection{$\bullet$\;\;\textbf{Sep. 2012 -- Jun. 2016, Harbin Institute of Technology, Harbin, China}}
\begin{spacing}{1.1}
  \quad\ Bachelor of Engineering, College of Civil Engineering
\end{spacing}
% \vspace{-2pt}
% \vspace{-12pt}
% \section{ \textcolor[rgb]{0,0.44,0.75}{\faChartBar\, 博士课题研究内容}}
\section{\faChartBar\, RESEARCH INTERESTS}
\begin{spacing}{1.05}
\begin{enumerate} [itemsep=3pt,topsep=2pt,parsep=0pt]
    \item Structural optimization
    \item Computational stochastic dynamics 
    \item Structural reliability analysis
    \item Reliability-based design optimization
    \item Reliability-based topology optimization
\end{enumerate}
\end{spacing}

\vspace{-12pt}
% \section{\textcolor[rgb]{0,0.44,0.75}{\faBook\;\; 学术成果}}
\section{\faBook\;\; PUBLICATIONS}

\begin{spacing}{1.1}
\subsection{\faFile*\,\; \textbf{Journal Papers}}
\vspace{10pt}
\begin{enumerate} [itemsep=3pt,topsep=2pt]

  \item Wan ZQ, Xie YQ, \underline{\textbf{Yang JS}}*, Lyu MZ, Ghafoori E. Accelerated stochastic model updating via decoupled multi-probability density evolution method with augmented change of probability measure. Mechanical Systems and Signal Processing 246: 113896. (IF=8.9, JCR Q1)

  \item Xiao R, Sushama L, Meguid M, Zayed T, \underline{\textbf{Yang JS}}*. (2026) Time-variant reliability analysis of corroded gas pipelines considering dynamic defect evolution and interaction. Reliability Engineering \& System Safety 267: 111899. (IF=11.0, JCR Q1)

  \item Chen JB*, Weng LL, \underline{\textbf{Yang JS}}. (2026) Dynamic-reliability-based design optimization by quantum-inspired algorithm and probability density evolution method combined with change of probability measure. Journal of Vibration Engineering 39(1):239-248. (In Chinese, EI)

  \item \underline{\textbf{Yang JS}}, Lyu MZ*, Chen JB, Xue JY. (2025) Structural design optimization under stochastic excitations considering first-passage probability constraint based on dimension-reduced probability density evolution equation. Reliability Engineering \& System Safety 264: 111378. (IF=11.0, JCR Q1)
  
  \item \underline{\textbf{Yang JS}}, Wan ZQ*, Jensen HA. (2025) An efficient strategy for information reuse in probability density evolution method considering large shift of distributions with multiple random variables. Probabilistic Engineering Mechanics 79: 103728. (IF=3.0, JCR Q1)
  
  \item Lyu MZ, \underline{\textbf{Yang JS}}*, Chen JB, Li J. (2025) High-efficient non-iterative reliability-based design optimization based on the design space virtually conditionalized reliability evaluation method. Reliability Engineering \& System Safety 254: 110646. (IF=9.4, JCR Q1)

  \item Chen G, Liu Y, Ma ZF, \underline{\textbf{Yang JS}}, Beer M, Chian SC*, Wei SJ*. (2025) Assessing extent of building damage following an earthquake: Case study of the 2023 Turkey-Syria doublet. npj Natural Hazards 2: 51. 

  \item \underline{\textbf{Yang JS}}, Chen JB*, Beer M. (2024) Seismic topology optimization considering first-passage probability by incorporating probability density evolution method and bi-directional evolutionary structural optimization. Engineering Structures 314: 118382. (IF=5.5, JCR Q1)

  \item Weng LL, Acevedo CH, \underline{\textbf{Yang JS}}, Valdebenito MA , Faes MGR , Chen JB*. (2024) An approximate decoupled reliability-based design optimization method for efficient design exploration of linear structures under random loads. Computer Methods in Applied Mechanics and Engineering 432: 117312. (IF=6.9, JCR Q1)

  \item Chen G*, \underline{\textbf{Yang JS}}, Wang RH, Li KQ, Liu Y, Beer M. (2024) Response to discussion of “Seismic damage analysis due to near-fault multipulse ground motion”.  Earthquake Engineering \& Structural Dynamics 53(2): 861–866. (IF=4.5, JCR Q1)

  \item Chen G*, \underline{\textbf{Yang JS}}, Liu Y, Kitahara T, Beer M. (2023) An energy-frequency parameter for earthquake ground motion intensity measure. Earthquake Engineering \& Structural Dynamics 52(2): 271-284. (IF=4.060, JCR Q1)

  \item Weng LL, \underline{\textbf{Yang JS}}, Chen JB*, Beer M. (2023) Structural design optimization under dynamic reliability constraints based on probability density evolution method and quantum-inspired optimization algorithm. Probabilistic Engineering Mechanics 74: 103494 (IF=2.6, JCR Q1)

  \item Chen G*, \underline{\textbf{Yang JS}}, Wang RH, Li KQ, Liu Y, Beer M. (2023) Seismic damage analysis subjected to near-fault multi-pulse ground motion. Earthquake Engineering \& Structural Dynamics 52(15): 5099-5116. (IF=4.5, JCR Q1, \textbf{Top cited and top viewed paper of the journal EESD})  
  
  \item Feng CX*, Faes MGR, Broggi M, Dang C, \underline{\textbf{Yang JS}}, Zheng ZB, Beer M, (2023). Application of interval field method to the stability analysis of slopes in presence of uncertainties. Computers and Geotechnics 153: 105060. (IF=5.218, JCR Q1)

  \item \underline{\textbf{Yang JS}}, Chen JB*, Beer M, Jensen HA. (2022) An efficient approach for dynamic-reliability-based topology optimization of braced frame structures with probability density evolution method. Advances in Engineering Software 173: 103196. (IF=4.255, JCR Q1)

  \item \underline{\textbf{Yang JS}}, Jensen HA, Chen JB*. (2022) Structural optimization under dynamic reliability constraints utilizing probability density evolution method and metamodels in augmented input space. Structural and Multidisciplinary Optimization 65(4): 107. (IF=4.542, JCR Q1)
  
  \item \underline{\textbf{Yang JS}}, Chen JB*, Jensen HA. (2022) Structural design optimization under dynamic reliability constraints based on the probability density evolution method and highly-efficient sensitivity analysis. Probabilistic Engineering Mechanics 68: 103205.
  (IF=3.350, JCR Q1)

  \item \underline{\textbf{Yang JS}}, Chen JB*. (2022) Structural optimization under dynamic reliability constraints by synthesizing change of probability measure and globally convergent method of moving asymptotes. Journal of Vibration Engineering 35(1):72-81. (In Chinese, EI, \textbf{High-influential paper of CNKI})

  \item Chen JB*, \underline{\textbf{Yang JS}}, Jensen HA. (2020) Structural optimization considering dynamic reliability constraints via probability density evolution method and change of probability measure. Structural and Multidisciplinary Optimization 62(5): 2499-2516. (IF=3.377, JCR Q1)  
\end{enumerate}
\end{spacing}

\vspace{-18pt}
\begin{spacing}{1.1}
% \subsection{\textcolor[rgb]{0,0.44,0.75}{\faUsers\ \textbf{会议论文}}}
\subsection{\faUsers\ \textbf{Conference Papers \& Presentations}}
\vspace{10pt}
\begin{enumerate}[itemsep=3pt,topsep=2pt]

  \item \underline{\textbf{Yang JS}}, Lyu MZ, Chen JB. (2024) Dynamic-reliability-based optimization of high-dimensional nonlinear structures under stochastic excitations using dimension-reduced probability density evolution equation. The 10th International Workshop on Reliable Engineering Computing (REC 2024). Beijing, China.

  \item \underline{\textbf{Yang JS}}, Ding YQ. (2024) Seismic-reliability-based topology optimization of frame structures considering stochastic ground motions. The 9th International Symposium on Reliability Engineering and Risk Management (ISRERM 2024). Hefei, China. (\textbf{Invited lecture of MS14})

  \item \underline{\textbf{Yang JS}}, Wan ZQ. (2024) Dynamic-reliability-based design optimization via probability density evolution method and efficient information reuse strategy. The 4th International Conference on Vulnerability and Risk Analysis and Management (ICVRAM2024) \& The 8th International Symposium on Uncertainty Modelling and Analysis (ISUMA2024). Shanghai, China. (\textbf{Invited lecture of MS06})
  
  \item \underline{\textbf{Yang JS}}, Chen JB. (2024) Reliability-based topology optimization of frame structures under seismic excitations by probability density evolution method. The 4th International Conference on Vulnerability and Risk Analysis and Management (ICVRAM2024) \& The 8th International Symposium on Uncertainty Modelling and Analysis (ISUMA2024). Shanghai, China.

  \item Weng LL, \underline{\textbf{Yang JS}}, Chen JB. (2023) Dynamic reliability-based design optimization using quantum-inspired algorithm and probability density evolution method assisted by change of probability measure strategy. The 14th International Conference on Applications of Statistics and Probability in Civil Engineering (ICASP14). Dublin, Ireland. (Full paper)

  \item \underline{\textbf{Yang JS}}, Chen JB, Jensen HA. (2022) A gradient-free approach for dynamic-reliability-based design optimization based on the probability density evolution method. The 13th International Conference on Structural Safety \& Reliability (ICOSSAR2021-2022). Shanghai, China. (Full paper)

  \item \underline{\textbf{Yang JS}}, Chen JB, Jensen HA. (2021) An efficient approach for dynamic-reliability-based design optimization using the probability density evolution method. The 14th World Congress on Structural and Multidisciplinary Optimization (WCSMO14). Boulder, USA.
  
  \item \underline{\textbf{Yang JS}}, Chen JB, Jensen HA. (2020) A PDEM-based method for structural optimization under dynamic reliability constraints. The 7th International Symposium on Reliability Engineering \& Risk Management (ISRERM2020). Beijing, China. (Full paper, \textbf{Student scholarship award})
  
  \item \underline{\textbf{Yang JS}}, Chen JB, Jensen HA. (2020) An efficient method for sensitivity estimation of failure probability based on the probability density evolution method. The 24th Annual Conference of HKSTAM 2020 \& The 16th Shanghai – Hong Kong Forum on Mechanics and Its Application. Hong Kong, China. (\textbf{Best student presentation award})
    
  \item \underline{\textbf{Yang JS}}, Chen JB, Jensen HA, Li J. (2019) Structural optimization considering dynamic reliability constraints via probability density evolution method. Proceedings of the 13th World Congress on Structural and Multidisciplinary Optimization (WCSMO13), 128-133. Beijing, China. (Full paper)
  
  \item \underline{\textbf{Yang JS}}, Chen JB, Jensen HA. (2019) An integrated method based on the probability density evolution method for dynamic reliability-based design optimization. The 32th KKHTCNN Symposium on Civil Engineering. Daejeon, South Korea. (Full paper)
\end{enumerate} 
\end{spacing}

\vspace{-18pt}
\begin{spacing}{1.1}
\subsection{\faFileCode\;\, \textbf{Code, Reports \& Preprints}}
\vspace{10pt}
\begin{enumerate}[itemsep=3pt,topsep=2pt]
  \item Chen JB, Lyu MZ, \underline{\textbf{Yang JS}}, Sun TT, Luo Y. (2023) MATLAB and Julia toolboxes for physically-driven DR-PDEE. (Code repository: \url{https://github.com/JCERSM/DR-PDEE-MATLAB} \& \url{https://github.com/JCERSM/DR-PDEE-Julia})
  
  \item \underline{\textbf{Yang JS}}, Chen JB. (2019) MATLAB toolbox for Probability Density Evolution Method (PDEM). (Code repository: \url{https://github.com/Tree-Yang/Process_Oriented_PDEM} \& \url{https://github.com/Tree-Yang/GF-discrepancy-based-point-selection})

\end{enumerate}
\end{spacing}
% \vspace{-18pt}
% \begin{spacing}{1.0}
% \subsection{\textcolor[rgb]{0,0.44,0.75}{\faFileSignature\, \textbf{在审论文}}}
% \subsection{\faFileSignature\, \heiti{在审论文}}
% \begin{enumerate}
%     \item \underline{\textbf{Yang JS}}, Chen JB, Beer M, Jensen HA. (2022) Dynamic-reliability-based topology optimization by probability density evolution method and bi-directional evolutionary structural optimization. Structural Safety. (在审;\, SCI: IF=5.712, JCR Q1.)
% \end{enumerate}
% \end{spacing}

% \section{\textcolor[rgb]{0,0.44,0.75}{\faAward\;\, 奖励荣誉}}

\section{\faArchive\;\, PROJECTS}
\begin{spacing}{1.1}
\begin{itemize} [itemsep=4pt,topsep=3pt]
  \item 2025.01-2027.12,  Young Scientists Fund (No. 52408212), National Natural Science Foundation of China (NSFC).
  
  \item 2024.01-2026.07, Special grants (No. 2023BSHTBZZ39), Postdoctoral Research Fund of Shaanxi Province, China.

  \item 2023.07-2026.07, Start-up funding for introduced talent (No. 1608724027), Xi'an University of Architecture and Technology, China.
\end{itemize}
\end{spacing}

\section{\faAward\;\, SCHOLARSHIPS \& AWARDS}
% increase linespacing [parsep=0.5ex]
 \begin{spacing}{1.1}
\begin{itemize} [itemsep=4pt,topsep=3pt]
  \item 2025, Top cited and top viewed article in Earthquake Engineering \& Structural Dynamics, Wiley.
  \item 2023, High-influential research paper, CNKI.
  \item 2021, Scholarship for studying abroad (12 months), Chinese Scholarship Council (CSC).
  \item 2020, Student scholarship award, the 7th International Symposium on Reliability Engineering and Risk Management (ISRERM2020).
  \item 2020, Best student presentation award, the 24th Annual Conference of Hong Kong Society of Theoretical and Applied Mechanics.
  \item 2020, First prize scholarship for civil engineering, Tongji University.
  \item 2016, Excellent doctoral (freshman) scholarship, Tongji University.
  \item 2015, First prize scholarship, Harbin Institute of Technology.
  \item 2013, National scholarship of China.
  \item 2013, Merit student, Harbin Institute of Technology.
\end{itemize}
 \end{spacing}

\vspace{-18pt}
% \section{ \textcolor[rgb]{0,0.44,0.75}{\faInfo\;\;\, 其他}}
% \section{\faInfo\;\;\, LANGUAGES \& SKILLS}
% increase linespacing [parsep=0.5ex]
% \begin{spacing}{1.1}
% \begin{itemize} [itemsep=2pt,topsep=2pt]
%   \item {\textbf{Languages}}：Chinese (First Language), English (Second Language, IELTS 6.5).
%   \item {\textbf{Programming}}：MATLAB, Julia, FORTRAN and Python for scientific computing.
%   \item {\textbf{Softwares}}: Abaqus, ANSYS and OpenSEES for finite element analysis.
% \end{itemize}
%  \end{spacing}
\section{\faInfo\;\;\, Academic Service}
\begin{spacing}{1.1}
\begin{itemize} [itemsep=2pt,topsep=2pt]
  \item Member of Specialized Committee on Safety of Urban Lifeline Engineering, China Association for Engineering Construction Standardization (CECS), 2023-
  \item Reviewer of international journals including: \\
  (1) Structural and Multidisciplinary Optimization\\
  (2) Computer Methods in Applied Mechanics and Engineering\\
  (3) Reliability Engineering \& System Safety\\
  (4) Mechanical Systems and Signal Processing\\
  (5) Engineering Structures\\
  (6) Probabilistic Engineering Mechanics\\
  (7) Chinese Journal of Aeronautics\\
  (8) ASCE-ASME Journal of Risk and Uncertainty in Engineering Systems, Part A: Civil Engineering\\
  (9) Journal of Building Engineering\\
  (10) Results in Engineering\\
  (11) World Earthquake Engineering (In Chinese)
  
  \item Member of International Society for Structural and Multidisciplinary Optimization (ISSMO).
\end{itemize}
 \end{spacing}
\end{document}
